\subsection{Mahler: ordered forms and boundary locality}
\label{sec:mahler}

The symmetric Mahler development connects a local differential-form calculation to a convex-body inequality. Exterior algebra describes the calculation, continuous alternating maps support its differentiation, and boundary integrals supply the quantity used in the geometric reduction. The selected endpoint, \src{OAI.SymmetricMahler.symmetric_mahler}, states
\[
  \frac{4^n}{n!}\leq |K|\,|K^\circ|
\]
for $n\geq1$ and a compact, convex, origin-symmetric set
$K\subseteq\R^n$ with nonempty interior. Here volume is ordinary Lebesgue
volume and the coordinate polar is literally
$K^\circ=\{p:\sum_i p_i v_i\leq1\text{ for every }v\in K\}$.
The accompanying argument uses holomorphic maps associated with finite intersections of strips~\citep{OAIMahler}. We follow the transgression identity through its regularity and locality requirements to the flux value consumed by the convex-body reduction.

\paragraph{The analytic inputs and the transgression identity.}
The generic analytic results take a structure named \src{MassHypotheses}.
Its fields require an open domain containing zero, holomorphic components
$f_j$, zero as their unique common zero in the domain, a homogeneous degree-$m$ leading
polynomial map $G$ without nonzero common zeros, an actual Euclidean Taylor
remainder $f-G=O(\norm{z}^{m+1})$, and compact closures of sublevels
with $0<R<1$ contained in the domain.
These fields record analytic inputs from which the mass and flux conclusions are derived. For the strip maps,
\src{euclideanSpecial_massHypotheses} constructs every field from positive
dimension, an injective real measurement matrix, and $m\geq2$.
The body theorem passes the resulting homogeneous flux value to the reduction lemma, leaving exactly the stated body hypotheses.\footnote{The endpoint and concrete hypothesis
construction are
\href{https://github.com/openai/math/blob/fd4aeeb2ee4fc729c18d98444fed42fd0529eeeb/lean/OAI/Analysis/Mahler/MainTheorem.lean\#L12-L22}{\src{Mahler/MainTheorem.lean}, lines 12--22}, and
\href{https://github.com/openai/math/blob/fd4aeeb2ee4fc729c18d98444fed42fd0529eeeb/lean/OAI/Analysis/Mahler/EuclideanSpecialMap.lean\#L67-L97}{\src{Mahler/EuclideanSpecialMap.lean}, lines 67--97}.}

For the transgression calculation in dimension $N\geq2$, write
$N=k+2$ with $k\geq0$, let $q(z)=z$, and set
$\tau(H)(z)=\sum_j|H_j(z)|^2$. The source convention is
$d^cu(z)[v]=-Du(z)[iv]/4$. On punctured space the degree-$m$
homogeneous map $G$ defines
\[
 \beta=d^c\log\tau(G)-m\,d^c\log\tau(q),\qquad
 \alpha_t=m\,d^c\log\tau(q)+t\beta.
\]
The deformation uses the ordered primitive
\[
  \eta_t=\alpha_t\wedge\beta\wedge(d\alpha_t)^k.
\]
The ambient differentiation theorem preserves both terms:
\[
  d\eta_t=\beta\wedge(d\alpha_t)^{k+1}
       -\alpha_t\wedge d\beta\wedge(d\alpha_t)^k.
\]
It follows from actual Fr\'echet differentiation, exterior product rules, and
$d^2\alpha_t=0$. Restriction to sphere tangents then kills the first term.
On sphere tangents, the remaining variation of $\alpha_t\wedge(d\alpha_t)^{k+1}$ contains $k+1$ placements of $d\beta$. Denote this remaining sum by $\operatorname{residual}_t$.
Its multiplicity gives the identity
\[
  \operatorname{residual}_t=-(k+1)d\eta_t
  \qquad\text{on sphere tangent tuples}.
\]
The two identities distinguish ambient differentiation from its restriction to the integration boundary. The sign and the factor $k+1$ identify the residual as an exact form there.

The disappearance of the first term has a concrete linear-algebra proof.
At $z\ne0$, the sphere tangent space has real dimension $2N-1$ and contains
the nonzero circle tangent $iz$. Homogeneity gives $\beta_z(iz)=0$ and
$(\alpha_t)_z(iz)=m/2$. Rotation invariance, together with this constant
circle evaluation, gives $d\alpha_t(iz,v)=0$.
Thus $\beta$ and $d\alpha_t$ are horizontal in this direction:
inserting $iz$ in any slot gives zero.\footnote{The normalization and
circle evaluations are in
\href{https://github.com/openai/math/blob/fd4aeeb2ee4fc729c18d98444fed42fd0529eeeb/lean/OAI/Analysis/Mahler/HomogeneousFlux.lean\#L12-L13}{the $d^c$ definition, lines 12--13}, and
\href{https://github.com/openai/math/blob/fd4aeeb2ee4fc729c18d98444fed42fd0529eeeb/lean/OAI/Analysis/Mahler/HomogeneousDeformation.lean\#L14-L101}{the deformation and circle identities, lines 14--101}.}
Their wedge has top degree on the sphere tangent space and therefore vanishes. The lemma
\src{horizontal_top_eq_zero} proves the general assertion explicitly.
An independent top-degree tuple determines a basis. Choose a nonzero basis
coordinate of the horizontal vector, replace the corresponding tuple entry
by that vector, and expand it in the basis. Alternation kills every term
except the selected coordinate times the original evaluation; cancelling the
nonzero coordinate proves that evaluation zero. A dependent tuple already
vanishes by alternation. The sphere theorem instantiates this argument with
the verified tangent dimension and $iz$.

\paragraph{Identifying two representations of the form.}
The development constructs $\eta_t$ in two ways suited to different operations. A literal shuffle form is bundled through finite basis expansion, supporting regularity proofs by evaluation. A continuous product of a two-form and wedge powers supports differentiation rules. Recursive slot types retain each factor's ordered
pair of arguments. Two equivalences flatten these slots to
$\operatorname{Fin}(2(k+1))$ along different routes. A small extensional proof
identifies the equivalences; injectivity of the projection from continuous
alternating maps to alternating maps then identifies the entire fields.
Equality of the whole fields transports the transgression into the analytic theorem expecting the other representation. The intermediate theorem
\src{homogeneousFlux_endpoints_of_transgression} has an explicit local-formula
premise, which \src{homogeneous_flux_endpoints} discharges.\footnote{See
\href{https://github.com/openai/math/blob/fd4aeeb2ee4fc729c18d98444fed42fd0529eeeb/lean/OAI/Analysis/Mahler/HomogeneousEndpoint.lean\#L12-L45}{the homogeneous endpoint, lines 12--45}.
The underlying tangent identity is proved in
\href{https://github.com/openai/math/blob/fd4aeeb2ee4fc729c18d98444fed42fd0529eeeb/lean/OAI/Analysis/Mahler/HomogeneousTransgression.lean\#L65-L97}{the tangent transgression theorem, lines 65--97}.}

The shuffle description exposes permutation signs and factor order for algebraic rewrites. The continuously bundled field has a normed-space structure, so the calculus library interface can express its regularity and derivative.
Finite basis expansion connects these obligations: smoothness of the finitely
many evaluated coefficients supplies smoothness of the bundled form.
The source obtains every finite regularity order on punctured space from the
analytic hypotheses, including the $C^2$ order used below.
Identifying the fields also identifies their derivatives as functions valued
in continuous alternating maps. Agreement on tangent tuples supplies the later integration comparison; differentiation uses equality of the form-valued fields. The proof therefore uses two distinct congruence principles at
appropriate stages: whole-field equality for differentiation and tangential
agreement for oriented density. The source uses each principle at the stage where its information suffices.

The continuous product also exposes a useful implementation choice.
\src{continuousWedge} starts with the algebraic shuffle product and proves
continuity in its vector arguments term by term. For continuity as a
function of the two forms, the source first embeds continuous alternating
maps into multilinear maps. Injectivity gives finite dimensionality of
the form space when the ambient real space is finite dimensional.
It then bundles the product as
\[
 \mathcal W:\mathcal A_\iota\longrightarrow_L
       (\mathcal A_\kappa\longrightarrow_L\mathcal A_{\iota\sqcup\kappa}),
\]
where $\mathcal A_\iota$ is the normed space of continuous alternating
maps indexed by $\iota$, and $\longrightarrow_L$ denotes a continuous
linear map over $\R$. The two applications of
\src{LinearMap.toContinuousLinearMap} use finite dimensionality to promote
algebraic linearity to continuity. The regularity of
$x\mapsto\mathcal W(a(x))(b(x))$ now follows from the ordinary calculus
of continuous linear maps and their application.\footnote{See
\href{https://github.com/openai/math/blob/fd4aeeb2ee4fc729c18d98444fed42fd0529eeeb/lean/OAI/Analysis/Mahler/ContinuousWedge.lean\#L19-L107}{\src{ContinuousWedge.lean}, lines 19--107}.}
The finite-dimensional setting makes this bridge economical: the bundled bilinear product supplies the calculus operations, and the shuffle representation supplies exact algebraic identities. Mathlib
already supplies the exterior derivative and its square-zero theorem
\citep{MathlibForms}; the project-specific work here connects the ordered
product, its regularity, and the integration interfaces.

\paragraph{From complex forms to real-coordinate flux.}
The integration interface requires another preservation theorem. The primitive is a complex-valued real alternating form, whereas coordinate Stokes takes real-valued forms on $\R^{2N}$. Let $L:\R^{2N}\to\mathbb C^N$ be the positively oriented orthonormal coordinate map used by the sphere flux. The source defines the real pullback
\[
 (\mathcal P\eta)_x(v_1,\ldots,v_q)
   =\operatorname{Re}\bigl(\eta_{Lx}(Lv_1,\ldots,Lv_q)\bigr)
\]
and proves $d(\mathcal P\eta)=\mathcal P(d\eta)$ wherever $\eta$ is differentiable. Taking real parts commutes with the finite alternating derivative sum; the chain rule handles the fixed linear coordinates. A separate density identity uses the chosen frame's volume $+1$, and a measure-preserving sphere isometry identifies the two flux integrals. Thus the exact-form identity reaches coordinate Stokes with its sign and normalization preserved.\footnote{See
\href{https://github.com/openai/math/blob/fd4aeeb2ee4fc729c18d98444fed42fd0529eeeb/lean/OAI/Analysis/Mahler/ExactFormCoordinates.lean\#L33-L96}{real-part projection and pullback naturality, lines 33--96}, and
\href{https://github.com/openai/math/blob/fd4aeeb2ee4fc729c18d98444fed42fd0529eeeb/lean/OAI/Analysis/Mahler/ExactSphereFluxBridge.lean\#L11-L67}{the oriented density and flux comparisons, lines 11--67}.}

\paragraph{Applying Stokes near the boundary.}
A second interface controls locality. The primitive is $C^2$ on punctured
space, while the flux argument needs regularity near the sphere.
\src{exists_closed_set_form_extension} takes a form $\omega$ that is $C^2$
on an open set $U$ containing a closed set $K$. Smooth separation gives a
function $\chi$ equal to one near $K$ and zero near $U^c$.
The field $\widetilde\omega=\chi\omega$ is globally $C^2$: use the product
rule inside $U$ and its eventual equality to zero outside $U$.
Near $K$ it is eventually equal to $\omega$, so their Fr\'echet derivatives,
and hence exterior derivatives, agree there.
For $K$ the sphere, coordinate Stokes applies to $d\widetilde\omega$ and
$d^2\widetilde\omega=0$ gives zero boundary flux. Tangent-density congruence
then transfers that conclusion to the original residual.\footnote{The
extension and exact-form flux theorem are
\href{https://github.com/openai/math/blob/fd4aeeb2ee4fc729c18d98444fed42fd0529eeeb/lean/OAI/Analysis/Mahler/ClosedBoundaryStokes.lean\#L18-L63}{\src{Mahler/ClosedBoundaryStokes.lean}, lines 18--63}.}

Zero residual flux makes the deformation's path flux constant by the
fundamental theorem of calculus. Comparing endpoints supplies $(\pi m)^N$;
the small-sphere limit and special strip-map reduction use that value in the
body inequality. Dimension one is handled by a separate circle calculation.
The mathematical devices are familiar. Mathlib already provides continuous
alternating maps, exterior differentiation, and $d^2=0$ under smoothness
assumptions~\citep{MathlibForms}. Kudryashov's earlier Lean formalization already transfers a divergence
theorem between integration theories and derives Cauchy--Goursat under
carefully chosen differentiability assumptions \citep{Kudryashov2022}.
Its rectangular domains give a direct comparison for integration
interfaces. The integration here combines equality of ordered and continuously bundled forms with a locality theorem supplying regularity near the boundary. This connects the analytic calculation to Stokes while accommodating the singular center. Section~\ref{sec:local-forms} illustrates why the ambient and tangential comparisons need different data.

\subsection{Crouzeix: compression, energy, and strict neighborhoods}
\label{sec:crouzeix}

The Crouzeix case follows a different analytic interface: boundary coefficient estimates feed a matrix-norm calculation with constant two.
For an $n$-by-$n$ complex matrix $A$, $m$-by-$m$ coefficients $B_k$, and
$n,m>0$, \src{OAI.DirectCrouzeix.complete_crouzeix} states
\[
 \norm{\sum_{k=0}^d A^k\otimes B_k}
 \leq 2\sup_{z\in W(A)}\norm{\sum_{k=0}^d z^kB_k}.
\]
Write $F(z)=\sum_{k=0}^d z^kB_k$; $F(A)$ denotes the amplified
evaluation on the left. The source uses the operator $2$-norm, literal matrix Kronecker products,
and $W(A)=\{\langle u,Au\rangle:\norm{u}=1\}$. Here $C^*$ denotes the
conjugate transpose, and the Hilbert--Schmidt (Frobenius) norm is
$\norm{C}_{\mathrm{HS}}^2=\sum_{i,j}|C_{ij}|^2$.
Its endpoint quantifies over arbitrary $A$ and coefficient matrices.
The accompanying argument~\citep{OAICrouzeix} is organized around boundary
analysis. The mathematical-frontiers review explains its ordered product
tests and positive-block closure in relation to earlier proof routes
\citep{GounderFrontiers2026}. Here the focus is the formal implementation:
zero-mode normalization, the finite dual witness, coefficient transport,
and construction of the strict neighborhoods used by the endpoint.

\paragraph{Closing the compressed energy estimate.}
The local interface \src{analytic_domain_core} accepts a polynomial sequence
$b_k$ with $b_0=1$, a continuous boundary parametrization, and a matrix resolvent field
$R(t)$. Its hypotheses give finite Faber expansions, lower and upper
coefficient-energy bounds, exact coefficients
$\widehat R(-k)=b_k(A)$, vanishing positive Fourier coefficients, and
$R+R^*\geq0$. These are the inputs of the local theorem.
\src{exterior_collar_bound} constructs them from exterior conformal data
and supporting halfplanes before applying the core. In particular,
the positive resolvent and coefficient bounds are constructed at this geometric stage before the public endpoint is assembled.

For the normalized core, assume $\norm{F(h(t))}\leq1$ on the boundary,
where $h$ is its parametrization. Let $T$ be the amplified polynomial
evaluation and $\gamma=\norm{T}>0$;
the zero case is separate. A unit singular pair satisfies
$Tx=\gamma y$ and $T^*y=\gamma x$. The supporting lemma obtains a norm
attainer by compactness, then proves the adjoint relation through a
squared-distance calculation. Reshaping vectors indexed by
$\operatorname{Fin}(n)\times\operatorname{Fin}(m)$ produces rectangular
matrices $X,Y$. All boundary integrals below use normalized Haar measure
$\mu$ on the period-one circle $\R/\mathbb{Z}$.
The proof forms
\[
 P=X^*RX,\quad Q=Y^*RY,\quad M=Y^*RX,\quad L=X^*RY,
 \qquad a=\int\norm{M(t)}_{\mathrm{HS}}^2\,d\mu(t).
\]
A cross-boundary pairing estimate makes $a=0$ incompatible with $\gamma>0$. The two singular relations then produce the diagonal bounds
\[
 \int\norm{P+P^*}_{\mathrm{HS}}^2\leq\frac{4a}{\gamma^2},
 \qquad
 \int\norm{Q+Q^*}_{\mathrm{HS}}^2\leq\frac{4a}{\gamma^2}.
\]
The normalization $b_0=1$ identifies the zero Fourier coefficients:
\[
 \widehat P(0)=X^*X,\qquad \widehat Q(0)=Y^*Y,\qquad
 \widehat L(0)^*=\widehat M(0).
\]
The first two are Hermitian; the last is the zero-frequency premise of
the cross-energy comparison. All four compressions inherit the vanishing
of the positive Fourier coefficients of $R$.
Fourier identities then give $a\leq\int\norm{M+L^*}_{\mathrm{HS}}^2$.
Compressing $R+R^*$ by the concatenated columns $[X\ Y]$ yields a positive
block matrix with these diagonal and off-diagonal blocks. For a positive
block matrix $D=\left(\begin{smallmatrix}U&Z^*\\Z&V\end{smallmatrix}\right)$
and $J=\operatorname{diag}(I,-I)$, positivity of $D$ and $JDJ$ gives
\[
 0\leq\operatorname{Re}\operatorname{tr}(DJDJ)
   =\norm U_{\mathrm{HS}}^2+\norm V_{\mathrm{HS}}^2
     -2\norm Z_{\mathrm{HS}}^2.
\]
Integrating this estimate gives
\[
 2\int\norm{M+L^*}_{\mathrm{HS}}^2
 \leq\int\norm{P+P^*}_{\mathrm{HS}}^2
       +\int\norm{Q+Q^*}_{\mathrm{HS}}^2.
\]
Combining the inequalities gives $2a\leq8a/\gamma^2$. Since $a>0$ and $\gamma>0$, multiplication by $\gamma^2$ and cancellation yield $\gamma^2\leq4$, hence $\gamma\leq2$.\footnote{The four compressions and closing cancellation
appear in
\href{https://github.com/openai/math/blob/fd4aeeb2ee4fc729c18d98444fed42fd0529eeeb/lean/OAI/Analysis/DirectCrouzeix/DomainCore.lean\#L15-L117}{\src{DirectCrouzeix/DomainCore.lean}, lines 15--117}.}

\paragraph{Where the sharp energy weights come from.}
Two ingredients fix the constants in this reduction. First,
\src{singular_diagonal_estimates} tests both ordered products $FG$ and $GF$, where $G$ is a test polynomial.
The forward and adjoint singular relations turn them into different diagonal
compressions, preserving the noncommutative order of matrix coefficients.
The dual step is explicit: from bounds on every finite test family $C_k$,
\src{weighted_coefficient_dual} substitutes
$C_k=w_k^{-1}\overline{P_k}$ in the project's bilinear entry pairing,
where the bar denotes entrywise complex conjugation.
Both pairing and test energy become
$s=\sum_k\norm{P_k}_{\mathrm{HS}}^2/w_k$, leaving $s^2\leq bs$ and hence
$s\leq b$. Nonnegative bounded partial sums subsequently give the infinite
sequence estimate. The infinite bound is thus obtained from explicit finite optimizing witnesses.

Second, the weights are derived by boundary analysis.
A Faber series splits into a nonnegative Fourier series and a remainder
with no nonnegative coefficients. Orthogonality gives its energy as
the sum of coefficient energy and remainder energy. The remainder is also
the action of a nonnegative kernel with both marginals one on the
nonconstant series. A weighted squared-norm inequality and Fubini prove
kernel contraction. Consequently $w_0=1$ and $w_k=2$ for $k>0$.
The resulting half-weighted coefficient tail fits the exact Fourier identity
\[
 \int\norm{P+P^*}_{\mathrm{HS}}^2
 =4\left(\norm{\widehat P(0)}_{\mathrm{HS}}^2
    +\frac12\sum_{k\geq1}\norm{\widehat P(-k)}_{\mathrm{HS}}^2\right).
\]
Its factor four at frequency zero uses Hermiticity of $\widehat P(0)$;
the negative coefficients contribute in paired positive and negative
frequencies.\footnote{See
\href{https://github.com/openai/math/blob/fd4aeeb2ee4fc729c18d98444fed42fd0529eeeb/lean/OAI/Analysis/DirectCrouzeix/BoundaryEnergy.lean\#L57-L88}{\src{BoundaryEnergy.lean}, lines 57--88}, and
\href{https://github.com/openai/math/blob/fd4aeeb2ee4fc729c18d98444fed42fd0529eeeb/lean/OAI/Analysis/DirectCrouzeix/PositiveBlocks.lean\#L22-L62}{\src{PositiveBlocks.lean}, lines 22--62}.}

A further interface supplies matrix-valued Faber expansions from scalar
polynomials. A scalar sequence with its $k$th polynomial of degree $k$ is
mapped into polynomials over the matrix coefficient ring by sending scalars
to scalar matrices. Its leading coefficients are units. The implementation
then applies Mathlib's degree-filtered spanning theorem for polynomial
sequences and extracts a finite matrix-coefficient witness from membership
in that span. This gives the exact finite expansion consumed by the core.
Evaluation and Kronecker-evaluation identities carry the finite witness to the two matrix calculations. The established polynomial basis theorem is
reused across the change of coefficient ring, while explicit transport lemmas
preserve the order of multiplication. Together with rectangular reshaping,
this explains how the scalar boundary analysis supplies estimates uniformly
in the matrix coefficient size, the defining feature of the complete bound.

\paragraph{Constructing the analytic inputs.}
The analytic inputs of the core are built before the final limiting step.
For an exterior parametrization $\psi(\zeta)=c\zeta+g(\zeta^{-1})$, the
source uses the algebra-valued generating function
\[
 \mathcal G_A(u)=
 (c-u^2g'(u))\bigl((c+ug(u))I-uA\bigr)^{-1}.
\]
Its $k$th Taylor coefficient is the Faber polynomial $b_k$ evaluated at
$A$. The proof makes this coefficient identity finite: a geometric
expansion of the inverse has a remainder $u^{k+1}H(u)$ with $H$ analytic
near zero, whose $k$th coefficient vanishes. Local equality transfers the Taylor coefficients. On the analytic disk supplied by the collar construction, evaluation at $u=e^{-2\pi it}$, with $t\in\R/\mathbb{Z}$, invokes the antianalytic Fourier formula, giving
the negative modes $b_k(A)$ and zero positive modes. Thus the orientation
and normalization consumed by \src{analytic_domain_core} arise from a
proved coefficient calculation.\footnote{See
\href{https://github.com/openai/math/blob/fd4aeeb2ee4fc729c18d98444fed42fd0529eeeb/lean/OAI/Analysis/DirectCrouzeix/FaberCalculus.lean\#L94-L163}{\src{FaberCalculus.lean}, lines 94--163}.}

Constructing $\psi$ requires a further change of representation. The source
inverts the exterior of a convex exponential sublevel about an interior
point, obtains a bounded simply connected domain, and constructs its
Riemann map. To obtain an analytic collar, it first extends that map
locally through analytic boundary charts. The gluing argument retains
more than analyticity: each local extension $E$ agrees with the original
map inside the domain, has nonzero derivative at its boundary point, and
satisfies
\[
 \norm{E(z)}<1\quad\Longleftrightarrow\quad z\in U.
\]
This side condition carries the geometric meaning of the unit circle
through continuation.

The local-to-global step is explicit in Lean. Extensions are available
on balls of radius $3r$, while the chosen cover uses balls of radius $r$.
When two smaller balls overlap, their larger intersection is convex and
meets $U$; the identity theorem extends their agreement from the interior
to that intersection. A classically chosen extension at each point is
therefore locally equal to any covering representative. This neighborhood
equality supplies analyticity and derivative agreement for the glued map.
Compactness and local invertibility then give an injective neighborhood
of $\overline U$. Its image contains the closed unit disk, so a thickening
contains a disk of radius $R>1$. The inverse on that disk is analytic,
and inversion back to the original plane gives the exterior Laurent
form.\footnote{See
\href{https://github.com/openai/math/blob/fd4aeeb2ee4fc729c18d98444fed42fd0529eeeb/lean/OAI/Analysis/DirectCrouzeix/Continuation.lean\#L31-L111}{compatibility and gluing of boundary germs},
\href{https://github.com/openai/math/blob/fd4aeeb2ee4fc729c18d98444fed42fd0529eeeb/lean/OAI/Analysis/DirectCrouzeix/ConformalCollar.lean\#L129-L238}{the inverse collar construction}, and
\href{https://github.com/openai/math/blob/fd4aeeb2ee4fc729c18d98444fed42fd0529eeeb/lean/OAI/Analysis/DirectCrouzeix/ExteriorMapping.lean\#L148-L214}{the exterior map for exponential sublevels}.}
The collar therefore retains analytic germs, the domain-side identity, and injectivity. These are the geometric facts needed to justify the resolvent and coefficient estimates used earlier.

\paragraph{Removing the auxiliary domain assumptions.}
The last step supplies strict-domain geometry for an arbitrary numerical range. For every compact,
nonempty convex $K\subseteq\mathbb C$, \src{exists_exponential_outer}
constructs a compact sublevel of a finite sum of exponentials of affine
support inequalities. A finite direction net controls approximation;
four coordinate directions ensure boundedness. The sublevel contains $K$
strictly and lies in any prescribed neighborhood of $K$.
In the endpoint, set
$M=\sup_{W(A)}\norm{F}$ and $\rho=M+\varepsilon/2$.
Compactness places such an outer sublevel inside the open set
$\{z:\norm{F(z)}<\rho\}$. The scaled domain estimate gives
$\norm{F(A)}\leq2M+\varepsilon$, and removing $\varepsilon>0$ gives the
desired bound. The auxiliary outer domain supplies smooth strict containment, leaving positive matrix sizes and arbitrary coefficient matrices in the endpoint.\footnote{See
\href{https://github.com/openai/math/blob/fd4aeeb2ee4fc729c18d98444fed42fd0529eeeb/lean/OAI/Analysis/DirectCrouzeix/OuterApproximation.lean\#L59-L122}{\src{DirectCrouzeix/OuterApproximation.lean}, lines 59--122}, and
\href{https://github.com/openai/math/blob/fd4aeeb2ee4fc729c18d98444fed42fd0529eeeb/lean/OAI/Analysis/DirectCrouzeix/CompleteBound.lean\#L65-L97}{\src{CompleteBound.lean}, lines 65--97}.}

The Crouzeix--Palencia proof centers the Cauchy transform
of conjugate holomorphic functions and obtains the complete constant
$1+\sqrt2$~\citep{CrouzeixPalencia2017}. Its established bound also supplies
the baseline in the companion numerical-range certificate study, which
develops finite enclosures, error budgets, and a rational software
prototype \citep{GounderNumerical2026}. The present case follows the
identities and geometric hypotheses that connect the formal proof's
local analytic core to its public matrix theorem.
